\documentclass[10pt,a4paper]{amsart}
\usepackage[margin=19mm]{geometry}
\usepackage{amsmath,amssymb,mathtools}
\usepackage[scr=rsfs,cal=euler]{mathalfa}
\usepackage{xfrac}
\usepackage[unicode,hidelinks,bookmarks=false]{hyperref}
\newcommand{\ie}{i.e.\ }
\newcommand{\cf}{cf.\ }
\newcommand{\resp}{respectively\ }
\newcommand{\Subsection}{\subsection}
\providecommand{\nobreakdash}{\nobreak\hbox{-}}
\setlength{\emergencystretch}{2em}
\multlinegap0pt
\usepackage{mathtools}
\usepackage{xcolor}
\usepackage[shortlabels]{enumitem}
\usepackage{amssymb}
\newtheorem{theorem}{Theorem}
\usepackage{cleveref}
\crefname{theorem}{Theorem}{Theorems}
\def\E{\mathbb E}
\def\R{\mathbb R}
\newcommand{\bfF}{\mathbf{F}}
\newcommand{\bff}{\mathbf{f}}
\newcommand{\bfp}{\mathbf{p}}
\newcommand{\bfx}{\mathbf{x}}
\newcommand{\diff}{\mathrm{d}}
\title{Convergence, design and training of continuous-time dropout as a random batch method}
\author{Antonio \'Alvarez-L\'opez, Mart\'in Hern\'andez}
\date{Source: arXiv:2510.13134v2 (CC BY 4.0)}
\begin{document}
\maketitle

\noindent\emph{Source and licence.} Convergence, design and training of continuous-time dropout as a random batch method. Version arXiv:2510.13134v2 (CC BY 4.0), distributed by arXiv under the Creative Commons Attribution 4.0 International licence, http://creativecommons.org/licenses/by/4.0/, as recorded in the arXiv licence metadata for this exact version and shown on the abstract page https://arxiv.org/abs/2510.13134v2. Full licence notice: \texttt{licenses/arxiv-2510.13134-CC-BY-4.0.txt}.\par\smallskip

\noindent\textit{Editorial scope.} Counted unit: the article's theorem thm:cost\_consistency (source0.tex lines 639--651). The article's complete original proof of it is delivered with it (source0.tex lines 1864--1942, one \texttt{proof} environment of 79 lines, which the article places away from the statement). No mathematics is written, repaired or re-derived: the statement and the proof are the article's own lines, in the article's own notation, and no retained line is substituted or re-flowed.  Retained uncounted as the source-local context the proof needs: lines 229--234; lines 288--298; lines 405--410; lines 473--479; lines 476--478.  Nothing was omitted from any retained span, so the package carries no omission record.  No retained line contains a citation, so the wrapper carries no bibliography and no bibliography entry.  No retained line contains a figure environment, an image-inclusion command or drawing code, so no image is delivered and the package declares no asset dependency.  Editorial formatting only, in addition: an AMS \texttt{amsart} preamble that restates the article's own declarations and macros used by the retained lines, namely the environments \texttt{theorem} on separate counters, together with the standard packages those lines need.  The retained environments print in this document as Theorem 1 in order of appearance, whereas the article prints the same items with the numbering of its own full document.  The complete native source of the article as distributed in the arXiv e-print for this exact version is bundled unchanged as \texttt{sources/187095/source0.tex}.
\begin{align}\tag{$\mathsf{FM}$}\label{eq:FM}
\begin{cases}
\dot\bfx_t&=\bfF(\bfx_t,\vartheta_t), \hspace{1cm}\text{for a.e. }t\in(0,T),\\
\bfx_0&\in\R^d,
\end{cases}
\end{align}

\begin{align}
0\le\lambda_{\bff_i,\theta,\mathsf B,\mathsf K_i}<\infty
&\qquad
\text{for all compact sets $\mathsf B\subset\R^d$ and
$\mathsf K_i\subset\Theta_i$,}
\tag{A1}\label{ass:lipthetaF}\\
0\le\lambda_{\nabla_\bfx\bff_i,\bfx,\mathsf K_i}<\infty
&\qquad
\text{for every compact set $\mathsf K_i\subset\Theta_i$.}
\tag{A2}\label{ass:lipxdxF}
\end{align}

\begin{align}\tag{$\mathsf{RM}$}\label{eq:RM}
\begin{cases}
\dot{\hat{\bfx}}_t =\hat{\bfF}_t(\hat{\bfx}_t,\vartheta_t), \hspace{1cm}&\text{for a.e. }t\in(0,T),\\
\hat{\bfx}_0=\bfx_0,
\end{cases}
\end{align}

\begin{theorem}\label{th:convergence}
For any $\bfx_0\in\R^d$ and $\vartheta\in L^\infty(0,T;\Theta)$, there is a finite constant
$\mathsf S(\bfx_0,\vartheta)\ge0$, independent of $h$, such that for every $n_s \ge 1$ (with $h=T/n_s$), the solutions $\bfx$ of \eqref{eq:FM} and $\hat{\bfx}$ of \eqref{eq:RM} satisfy
\begin{equation}\label{eq:bound.th:convergence}
\max_{t\in[0,T]}\E_{\omega}\left[\|\bfx_t-\hat{\bfx}_t\|^2\right]\leq  \mathsf{S}(\bfx_0,\vartheta) \, h.  
\end{equation}
\end{theorem}

\begin{equation}
\max_{t\in[0,T]}\E_{\omega}\left[\|\bfx_t-\hat{\bfx}_t\|^2\right]\leq  \mathsf{S}(\bfx_0,\vartheta) \, h.  
\end{equation}

\begin{theorem}\label{thm:cost_consistency}
There exist constants $C_{\mathcal U}^{\rm str},C_{\mathcal U}^{\rm wk}>0$, independent of $h$, such that:
\begin{enumerate}
    \item Without additional assumptions,
    \begin{equation}\label{eq:strong_j_bound}
    \sup_{\vartheta\in\mathcal U} \mathbb E_\omega\left[ |\hat J_h(\vartheta,\omega)-J(\vartheta)|^2 \right] \le C_{\mathcal U}^{\rm str}h.
    \end{equation}
   \item If \Cref{ass:lipxdxF} holds, then
    \begin{equation}\label{eq:weak_j_bound}
    \sup_{\vartheta\in\mathcal U} |\overline J_h(\vartheta)-J(\vartheta)| \le C_{\mathcal U}^{\rm wk}h.
    \end{equation}
\end{enumerate}
\end{theorem}

\begin{proof}[Proof of \Cref{thm:cost_consistency}]
\textbf{Strong bound \eqref{eq:strong_j_bound}:} By the uniform version of \Cref{th:convergence}, for each $m\in[n_d]$
there exists $C_{m,\mathcal U}>0$ such that
\[
\sup_{\vartheta\in\mathcal U}
\max_{t\in[0,T]}
\E_\omega[
\|\hat\bfx_{m,\vartheta,t}-\bfx_{m,\vartheta,t}\|^2]
\le C_{m,\mathcal U}h.
\]

The control penalty cancels out when taking the difference of the objectives. Applying Minkowski's inequality in $L^2(\Omega)$, and using the fact that $\ell_m, g_m \in {C}^{1,1}_{\rm loc}(\R^d; \R_{\geq0})$ imply that $\ell_m$ and $g_m$ are locally Lipschitz for every $m\in[n_d]$, we have
\begin{align*}
\|\hat J_h(\vartheta,\cdot) - J(\vartheta)\|_{L^2(\Omega)} \le& \sum_{m\in[n_d]} \bigg[ \beta\lambda_{\ell_m,B_R} \int_0^T \|\hat\bfx_{m,t}-\bfx_{m,t}\|_{L^2(\Omega)} \diff t \\
& + \lambda_{g_m,B_R} \|\hat\bfx_{m,T}-\bfx_{m,T}\|_{L^2(\Omega)} \bigg],
\end{align*}
where $B_R \subset \R^d$ is a compact ball containing all trajectories uniformly. Substituting \eqref{eq:bound.th:convergence} into this expression gives $\|\hat J_h(\vartheta,\cdot) - J(\vartheta)\|_{L^2(\Omega)} \le \widetilde{C}_{\mathcal U}\sqrt{h}$. Squaring both sides yields \eqref{eq:strong_j_bound}.

\textbf{Weak bound \eqref{eq:weak_j_bound}:} Fix a deterministic control $\vartheta \in \mathcal U$ and let $e_{m,t} \coloneqq \hat\bfx_{m,t} - \bfx_{m,t}$. By Taylor expansion of the locally smooth losses,
\begin{align*}
g_m(\hat\bfx_{m,T}) - g_m(\bfx_{m,T}) &= \langle \nabla g_m(\bfx_{m,T}), e_{m,T} \rangle + \mathcal{O}(\|e_{m,T}\|^2), \\
\ell_m(\hat\bfx_{m,t}) - \ell_m(\bfx_{m,t}) &= \langle \nabla \ell_m(\bfx_{m,t}), e_{m,t} \rangle + \mathcal{O}(\|e_{m,t}\|^2).
\end{align*}
Since $\sup_{t} \mathbb E_\omega[\|e_{m,t}\|^2] \le C h$ and the trajectories remain in a compact ball $B_R$, the expected quadratic remainders are uniformly bounded by $\mathcal{O}(h)$. To control the expected linear terms, we introduce the deterministic adjoint state $\bfp_{m,t}$ associated with the full model, which solves
\begin{equation*}
\begin{cases}
\dot\bfp_{m,t} = -\nabla_\bfx\bfF(\bfx_{m,t},\vartheta_t)^\top\bfp_{m,t} - \beta\nabla\ell_m(\bfx_{m,t}), \hspace{1cm} \text{for a.e. } t\in (0,T),\\
\bfp_{m,T}=\nabla g_m(\bfx_{m,T}).
\end{cases}
\end{equation*}
Defining the constants $M_{g_m,R} \coloneqq \sup_{\bfx\in B_R} \|\nabla g_m(\bfx)\|$ and $M_{\ell_m,R} \coloneqq \sup_{\bfx\in B_R} \|\nabla \ell_m(\bfx)\|$, Grönwall's inequality guarantees that this adjoint state is uniformly bounded over $\mathcal{U}$:
\begin{equation*}
\sup_{\vartheta\in\mathcal U}\|\bfp_{m,\vartheta}\|_{L^\infty(0,T)} \le e^{\lambda_{\bfF,\bfx,\mathsf K}T}\left(M_{g_m,R} + \beta T M_{\ell_m,R}\right).
\end{equation*}
Since $e_{m,0}=0$, $\bfp_{m,T}=\nabla g_m(\bfx_{m,T})$, and $\beta\nabla\ell_m(\bfx_{m,t})
=
-\dot\bfp_{m,t}
-
\nabla_\bfx\bfF(\bfx_{m,t},\vartheta_t)^\top\bfp_{m,t},$ integration by parts gives
\begin{align*}
&\left\langle\nabla g_m(\bfx_{m,T}),e_{m,T}\right\rangle
+\beta\int_0^T
\left\langle\nabla\ell_m(\bfx_{m,t}),e_{m,t}\right\rangle\diff t
\int_0^T
\left\langle
\bfp_{m,t},
\dot e_{m,t}
-
\nabla_\bfx\bfF(\bfx_{m,t},\vartheta_t)e_{m,t}
\right\rangle\diff t.
\end{align*}
Summing over $m$ and taking expectations gives the total expected linear
contribution. To evaluate its integrand, we expand the randomized vector
field around the deterministic state:
\begin{align*}
\dot e_{m,t} - \nabla_\bfx\bfF(\bfx_{m,t},\vartheta_t)e_{m,t} &= \hat\bfF_t(\bfx_{m,t},\vartheta_t) - \bfF(\bfx_{m,t},\vartheta_t) \\
&\quad + \underbrace{\left( \nabla_\bfx\hat\bfF_t(\bfx_{m,t},\vartheta_t) - \nabla_\bfx\bfF(\bfx_{m,t},\vartheta_t) \right)}_{\coloneqq G_{m,t}}e_{m,t} + \mathcal{R}_{m,t}.
\end{align*}
By \Cref{ass:lipxdxF} on the compact parameter sets
$\mathsf K_i=\operatorname{proj}_i\mathsf K$, the Taylor remainder
satisfies uniformly
\[
\|\mathcal R_{m,t}\|
\le C_{\mathcal U}\|e_{m,t}\|^2.
\]
Thus, its expectation yields a $\mathcal{O}(h)$ contribution. 

The first term evaluates the difference of the fields at the deterministic state $\bfx_{m,t}$; by unbiasedness, its expectation is identically zero. To control the critical cross-term $G_{m,t}e_{m,t}$, we introduce the filtration $\mathcal F_k \coloneqq \sigma(\omega_1,\ldots,\omega_k)$. Here the control $\vartheta$ is deterministic and the current batch is independent of the past filtration. On any interval $t \in [t_{k-1}, t_k)$, we decompose the error as $e_{m,t} = e_{m,t_{k-1}} + (e_{m,t} - e_{m,t_{k-1}})$. Conditioning on the past yields
\begin{equation*}
\mathbb E_\omega\left[ G_{m,t} e_{m,t_{k-1}} \right] = \mathbb E_\omega\left[ \mathbb E_\omega\left[ G_{m,t} \mid \mathcal F_{k-1} \right] e_{m,t_{k-1}} \right] = 0,
\end{equation*}
because $\omega_k$ is independent of $\mathcal F_{k-1}$ and $\mathbb E_\omega[\nabla_\bfx\hat\bfF_t(\bfx,\theta)] = \nabla_\bfx\bfF(\bfx,\theta)$. The remaining part involves the increment $(e_{m,t} - e_{m,t_{k-1}})$, which satisfies
\begin{equation*}
\|e_{m,t} - e_{m,t_{k-1}}\| \le \|\hat\bfx_{m,t}-\hat\bfx_{m,t_{k-1}}\| + \|\bfx_{m,t}-\bfx_{m,t_{k-1}}\| \le \int_{t_{k-1}}^t \left(\|\hat\bfF_s\| + \|\bfF\| \right) \diff s \le C h,
\end{equation*}
uniformly bounded because the vector fields are uniformly bounded on $B_R \times \mathsf K$. Since both the adjoint state $\bfp_{m,t}$ and the gradient difference $G_{m,t}$ are uniformly bounded, the inner product with this increment contributes at most $\mathcal{O}(h)$ inside the integral. 

All constants above depend only on $\mathsf K$, the data, the losses, $T$, and the constants in \Cref{ass:lipxdxF}, and are therefore uniform over $\vartheta\in\mathcal U$. Summing over $m$ and taking the supremum over $\mathcal U$ proves \eqref{eq:weak_j_bound}.
\end{proof}

\end{document}
